\section{Introduction}
\label{sec:introduction}

Retrieval-augmented generation (RAG) gives a language model access to
information beyond its parameters by conditioning generation on external
passages~\citep{rag_survey}. A fixed retrieval policy is nevertheless poorly
matched to heterogeneous questions. Retrieval can add latency and distracting
evidence when the model already knows a simple fact, yet a single query is
often insufficient for a multi-hop question whose decisive entity becomes
known only after an intermediate step~\citep{ircot,ren2023}. Adaptive RAG
therefore needs to decide not only \emph{what} to retrieve, but also
\emph{when} the current reasoning state warrants retrieval.

One appealing solution is to use the language model's own token distribution.
Uncertainty-triggered methods can act during generation without a separately
trained routing model~\citep{flare,dragin,targ}. In practice, however,
uncertainty is not synonymous with factual error. A model can distribute
probability among several harmless aliases or formatting variants while
remaining factually correct; it can also assign almost all probability to a
fluent hallucination. Standard entropy combines these cases because it records
both competition among high-probability candidates and dispersion into the
large low-probability tail.

This paper asks whether a more selective statistic of the token distribution
can provide a better confidence signal for adaptive retrieval. We introduce
\tailleak{}, which measures how much probability escapes the model's
own small set of leading candidates. If \(m_t\) is the probability mass of the
top-\(k\) tokens at generation step \(t\), the answer-level score is
\[
    \tl(a)=-\sum_{t\in a}\log m_t.
\]
The score is near zero when a small candidate set retains almost all
probability and increases monotonically as mass moves into the tail. Unlike
full entropy, it discards how probability is divided inside the leading set
and how the remaining mass is arranged within the tail. It is computed from
the logits already produced during greedy decoding, without extra generations,
labels, or router training.

We evaluate this signal inside \method{}, short for Self-certainty Evaluated
Evidence-gated Knowledge-retrieval. This dual-path framework for open-domain
question answering first
produces a direct answer and also decomposes the original question into
ordered, executable sub-questions. Each sub-question is tentatively answered
without retrieval; \tailleak{} then decides whether to retrieve and re-answer
that step. The decomposed steps are synthesized into a second candidate, and a
transparent late rule selects between the direct and decomposition candidates.
This design makes decomposition reversible: it can expose bridge entities for
retrieval without discarding the original-question answer.

Our evaluation uses four development-set traces spanning three multi-hop
datasets---\hotpot{}, \twowiki{}, and \musique{}---and the single-hop
\popqa{} benchmark. Among six representative confidence families,
\tailleak{} gives the highest macro AUROC for detecting wrong direct answers
(\(0.749\), versus \(0.726\) for maximum token entropy). In the tested
end-to-end configuration, the \tailleak{} system obtains \(41.88/48.90\)
macro EM/F1, compared with \(40.52/47.61\) when maximum token entropy is used.
The accuracy-first retrieval gate is effectively tied with always retrieving
while skipping \(6.16\%\) of decomposed steps; more aggressive thresholds show
a gradual accuracy--retrieval trade-off. The results do not imply that the
score is a correctness test: confidently wrong answers remain, and final
routing contributes only a small net gain.

Our contributions are:
\begin{itemize}
    \item We derive \tailleak{}, a simple, monotonic, training-free statistic
    that isolates probability mass leaving the leading token set and
    accumulates it across an answer.
    \item We instantiate the score as a sub-question retrieval gate in
    \method{}, a reversible dual-path framework for decomposition-based RAG.
    \item We compare representative confidence signals, retrieval behavior,
    and end-to-end QA across four datasets, while identifying the boundary
    created by confident errors and imperfect path routing.
\end{itemize}
